% !TEX root = ../nc-quadric-classification.tex
% 2026-11-21 / v1.40
\section{五頂点窓と noncubic 変形の全 Artin 認識}
\label{sec:noncubic-local-tails}

\paragraph{局所環で成立することと、一様な開で残ること.}
本節では系~\ref{cor:noncubic-etale-morita-root} の有限型近傍の
\emph{指定された noncubic 点における局所環}を取る。
その同じ局所環上で、全ての前向き Hom、全軌道環、canonical
Grothendieck 導来圏と Noetherian heart を構成する。
極小解消の Betti 数を族の中で固定せず、五頂点の非極小 bar 解消を使う。
さらに二次双対の中心・反中心を使って同じ根の中央 AS4 被覆を局所環上に構成し、
全 noncubic 点で任意の Artin 変形を相対 $\mathfrak H_2$ として認識する。
有限型の一つの開近傍上での全次数一様性と全幾何 fiber の $\HH$ 所属は別問題であり、
本節だけから noncubic 点を含む $\mathfrak H_2$ の開 chart は得られない。

\subsection{極小 Betti 数に依存しない bar 解消}

以下で $C$ は正の向きの $R$-線形 $\mathbb Z$ 代数、
$P_i=e_iC$、$S_i$ は頂点 $i$ に $R$ を置く加群とする。
積は $C_{ij}\otimes C_{jk}\to C_{ik}$ と書く。

\begin{lemma}[五頂点 bar 解消の持上げ]
\label{lem:five-window-bar-lifting}
次のいずれかを仮定する。
\begin{enumerate}
\item $R$ は可換局所環であり、全 $C_{ij}$ は有限自由である。
  以下の fiber 条件を剰余体上で仮定する。
\item $R$ は任意の可換環であり、全 $C_{ij}$ は有限 projective である。
  以下の fiber 条件を全ての幾何点上で仮定する。
\end{enumerate}
fiber 条件は、各 $S_i$ が長さ三の有限生成 projective 解消を持ち、
その項に現れる表現可能加群が $P_i,\ldots,P_{i+4}$ に限られ、
最後の項が $P_{i+4}$ 一個であり、さらに
\[
 \operatorname{Ext}_{C_{\bar x}}^q(S_i,P_j)=0
 \quad\text{unless }(q,j)=(3,i+4),\qquad
 \operatorname{Ext}_{C_{\bar x}}^3(S_i,P_{i+4})\simeq\kappa(\bar x)
 \tag{BW-fiber}\label{eq:five-window-fiber-dual}
\]
となることである。最後の項以外の極小 Betti 数は一定と仮定しない。

このとき $S_i$ は $C$ 上 perfect で、projective dimension は高々三である。
解消と AS 双対条件は任意の $R$-基底変更後にも成立し、
\[
 \mathbf R\operatorname{Hom}_{C}(S_i,P_j)\simeq
 \begin{cases}\mathcal L_i[-3],&j=i+4,\\0,&j\ne i+4\end{cases}
 \tag{BW-dual}\label{eq:five-window-relative-dual}
\]
である。ここで $\mathcal L_i$ は可逆 $R$-加群である。
\end{lemma}

\begin{proof}
全証明は相対モジュライ論文 \cite[\texttt{noncubic_local_tails}]{QuadricModuli2026} に置く。
\end{proof}

\begin{remark}[二次曲面の固定階数の五頂点複体]
\label{rem:five-window-fixed-ranks}
$\operatorname{rank} C_{i,i+r}=1,2,4,6,9$（$r=0,1,2,3,4$）の場合、
\eqref{eq:five-window-bar} は局所的に
\begin{align*}
0\longrightarrow P_{i+4}^{16}&\longrightarrow
 P_{i+3}^{8}\oplus P_{i+4}^{48}
 \longrightarrow P_{i+2}^{4}\oplus P_{i+3}^{16}\oplus P_{i+4}^{40}\\
 &\longrightarrow P_{i+1}^{2}\oplus P_{i+2}^{4}
       \oplus P_{i+3}^{6}\oplus P_{i+4}^{9}
 \longrightarrow P_i\longrightarrow S_i\longrightarrow0
\end{align*}
となる。最終添字 $i+4$ の係数複体の階数は、homological degree
$1,2,3,4$ の順に $9,40,48,16$ である。
極小形が cubic 型から noncubic 型へ変わっても、この bar 複体の階数は変わらない。
本節の方法が要求するのはこの固定階数の複体の完全性である。
\end{remark}

\begin{proposition}[五頂点解消による canonical tails]
\label{prop:five-window-canonical-tails}
前補題の仮定に加え、$C$ が距離高々二の成分で生成されるとする。
$\mathcal O_i=\pi P_i\in\operatorname{QGr}C$ と置くと、
\[
 \operatorname{Ext}^q(\mathcal O_i,\mathcal O_j)\simeq
 \begin{cases}
 C_{ji},&q=0,\\
 \operatorname{Hom}_R(C_{i,j-4},\mathcal L_{j-4}),&q=2,\ j\ge i+4,\\
 0,&\text{otherwise}
 \end{cases}
 \tag{BW-coh}\label{eq:five-window-tails-cohomology}
\]
である。$\mathbf R\omega$ の cohomological dimension は高々二であり、
任意の連続四項の直和 $T_a=\bigoplus_{i=a}^{a+3}\mathcal O_i$ は
$D(\operatorname{QGr}C)$ の compact tilting generator となる。
従って
\[
 D(\operatorname{QGr}C)\simeq D(\operatorname{Mod}\operatorname{End}(T_a)),
 \qquad D(\operatorname{QGr}C)^c\simeq\operatorname{Perf}\operatorname{End}(T_a).
 \tag{BW-Morita}\label{eq:five-window-tails-morita}
\]
全構成は任意の導来基底変更と整合する。
Noetherian 性はこの命題の仮定にも結論にも含めない。
\end{proposition}

\begin{proof}
全証明は相対モジュライ論文 \cite[\texttt{noncubic_local_tails}]{QuadricModuli2026} に置く。
\end{proof}

\subsection{良い noncubic 根の局所環上での全軌道と heart}

\begin{theorem}[noncubic 局所環における全次数 tails 認識]
\label{thm:noncubic-local-tails-effectivity}
系~\ref{cor:noncubic-etale-morita-root} の $(U,u)$ と Morita 根 $\mathsf s$ を取り、
$R=\mathcal O_{U,u}$ とする。
より一般に、閉点を $u$ へ送る Noetherian 局所 $\mathcal O_{U,u}$-代数を $R$ としてもよい。
$\Lambda=\Lambda_U\otimes R$ とし、
\[
 E_{2a}=\mathsf s^aF_0,\quad E_{2a+1}=\mathsf s^aF_1,
 \quad C_{ij}=\operatorname{Hom}_{\Lambda}(E_{-j},E_{-i})\quad(i\le j)
\]
と置く。$C_{ij}=0$（$i>j$）とし、積は射の合成とする。
次が成立する。
\begin{enumerate}
\item 全 $i\le j$ で $\mathbf R\operatorname{Hom}(E_i,E_j)$ は有限自由 $R$-加群の
  次数零成分に集中する。階数は閉 fiber の normal helix と等しく、
  偶数差 $2n$ で $(n+1)^2$、奇数差 $2n+1$ で $(n+1)(n+2)$ である。
  この同定と積は任意の導来基底変更と整合する。
\item 全軌道環
  $B=\bigoplus_{n\ge0}\operatorname{Hom}(F_0,\mathsf s^nF_0)$ は
  一次生成・二次・相対 Koszul で、全 $B_n$ は階数 $(n+1)^2$ の自由加群である。
  $B$ は左右 Noetherian であり、任意の有限型可換 $R$-代数への変更後も同様である。
\item 任意の可換 $R$-代数 $R'$ について canonical な比較
  \[
  D(\operatorname{QGr}B_{R'})\simeq D(\operatorname{Mod}\Lambda_{R'}),
  \qquad D(\operatorname{QGr}B_{R'})^c\simeq\operatorname{Perf}\Lambda_{R'}
  \]
  があり、$\pi B_{R'}$ を $F_0$ に、次数移動 $(1)$ を $\mathsf s$ に送る。
  特に compact 圏の canonical Serre 関手について $(1)^2\simeq S^{-1}[2]$ である。
\item 有限型可換 $R$-代数 $R'$ 上で $\mathcal C_{R'}=\operatorname{qgr}B_{R'}$ は
  Noetherian な相対 Abel 圏の系をなし、module-object 基底変更と fppf descent を持つ。
  coherent 対象間の Hom は有限で平坦基底変更と可換する。
  $\pi B_{R'}$ は平坦かつ exceptional で、$(\pi B_{R'},(1))$ は AZ ample である。
  左右の全 $\chi_q$ が成立し、canonical compact 圏は $\operatorname{Perf}\Lambda_{R'}$ である。
\end{enumerate}
この定理だけには同じ根の AS4 被覆を含めない。標準局所基底では、
後の定理~\ref{thm:noncubic-local-central-hull} がその被覆を与える。
また $R$ は一般に元の局所有限型 $k$-site の対象ではないため、
以上をそのまま $U$ の開近傍上の $\mathfrak H_2$ 族と呼ばない。
\end{theorem}

\begin{proof}
全証明は相対モジュライ論文 \cite[\texttt{noncubic_local_tails}]{QuadricModuli2026} に置く。
\end{proof}

\begin{remark}[局所定理と有限型一様化の役割]
有限型近傍上の各軌道 Hom は perfect であり、固定した有限個については開性を使える。
本定理は、その無限個の開条件の共通部分が開であるとは述べない。
AS4 被覆の延長が得られても、有限型近傍で使う全軌道の一様平坦性・表示比較と
全幾何 fiber の rank-one torsion-free 条件は別途確認する。
一方、同じ局所環上では全次数の bar 解消と Noetherian heart を既に得たので、
以後の局所 AS4 構成で canonical compact 圏を独立の未証明入力に戻す必要はない。
有限型一様化は定理~\ref{thm:noncubic-finite-type-effectivity} で、
全幾何 fiber の有限 foundation 認識を新しい入力として証明する。
\end{remark}


\subsection{中央被覆の局所延長と全 Artin 認識}

\begin{lemma}[cubic 根の canonical hull の scalar normalizer]
\label{lem:cubic-compatible-scalar-normalizer}
体 \(K\) 上の geometric cubic \(\mathbb Z\)-代数に実際の二周期 \(\beta\) があり、
対応する tails 自己同値 \(s\) が \(s^2\simeq S^{-1}[2]\) を満たすとする。
その偶数環 \(B\) の canonical hull の双対線 \(K\phi\subset B^!_2\) は、
ある \(\lambda\in\{1,-1\}\) に対して
\[
 \phi x=\lambda x\phi\qquad(x\in B^!_1)
\]
を満たす。ここでは最初から \(\lambda=1\) とは主張しない。
\end{lemma}
\begin{proof}
全証明は相対モジュライ論文 \cite[\texttt{noncubic_local_tails}]{QuadricModuli2026} に置く。
\end{proof}

\begin{theorem}[局所基底上の noncubic 横断中央被覆]
\label{thm:noncubic-local-central-hull}
\((R,\mathfrak m,k)\) を標数零の体上の Noetherian 正則局所環、\(K=\operatorname{Frac}R\)
とする。二次環 \(B=T_R(V)/(J)\) について次を仮定する。
\begin{enumerate}
\item \(V,J\) は階数四、七の自由加群で、\(J\subset V^{\otimes2}\) は直和因子である。
全ての \(B_n\) は階数 \((n+1)^2\) の有限自由加群であり、\(B\) は相対 Koszul である。
特に \(E=B^!\) の全次数は有限自由、任意の底変更と可換し、その階数は
\(1,4,7,8,8,\ldots\) である。
\item 閉 fiber \(B_k\) は定理~\ref{thm:noncubic-central-hull} の良い根の
noncubic 軌道環である。
\item generic fiber \(B_K\) は geometric cubic helix を二項移す実際の根の偶数環であり、
その根の平方は反 Serre 関手に同型である。
\end{enumerate}
このとき \(B\) は相対 quadratic Koszul AS4 環 \(A\) と中央正則二次元 \(g\) により
\[
 B\cong A/(g)
\]
と表示される。全次数の有限自由性、左右および双加群 Koszul 解消、中央正則切断は
任意の基底変更と両立する。\(B\) が左右 Noetherian なら \(A\) もそうである。
この定理は有限型開近傍上の一様な Noetherian 性をそれ自体では主張しない。
\end{theorem}

\begin{proof}
全証明は相対モジュライ論文 \cite[\texttt{noncubic_local_tails}]{QuadricModuli2026} に置く。
\end{proof}

\begin{remark}[Artin 結論と有限型一様性]
この定理の入力を有限 foundation の局所環上で確認できれば、
任意の Artin 厚みへの引戻しは、同じ良い根の AS4 被覆を与える。
しかし局所環上の全次数平坦性・Noetherian 性の証明だけから、
単一の有限型開近傍上で全次数の認識が成立すると推論しない。
有限型近傍の全 orbit、canonical tails と普遍 Noetherian 性はそれぞれ別の入力である。
\end{remark}

\begin{corollary}[全 noncubic 点の Artin 認識]
\label{cor:noncubic-all-artin-h2}
標数零の代数閉体上の任意の noncubic \(\mathcal H_2\) 対象
\((\mathcal C,\mathcal O)\) を取る。
その任意の flat pointed Abel Artin 変形は、指定した良い根の持上げの全軌道環に対する
中央正則二次 AS4 被覆を持ち、全有限型基底変更を備えた相対 \(\mathfrak H_2\) 族へ昇格する。
従って、局所 Artin 基底上で相対 \(\mathfrak H_2\) 変形から pointed Abel 変形への忘却は
変形 groupoid の同値である。
根対が異なる点では接次元三・無限小慣性零、根対が一致する点では接次元四・無限小慣性一次元であり、
いずれも変形は無障害である。
\end{corollary}
\begin{proof}
全証明は相対モジュライ論文 \cite[\texttt{noncubic_local_tails}]{QuadricModuli2026} に置く。
\end{proof}

\begin{remark}[形式的完了と代数性の区別]
全 cubic 点における既存の形式認識と前系を合わせると、全 \(\mathcal H_2\) 幾何点で
\(\mathfrak H_2\) と pointed Abel 変形の Artin groupoid が一致する。
この結論から内在的スタック全体の有限型代数性はまだ従わない。
局所環上で同時に成り立つ全軌道・Noetherian 条件を、単一の有限型 étale 近傍へ
一様に広げる段階と、その相対認識を別に閉じる必要がある。
\end{remark}

\subsection{中央二次元の全次数正則性の有限型一様化}

\begin{lemma}[有限な二次商から中央正則線を一様化する]
\label{lem:spread-central-regular-line}
$R$ を可換 Noetherian 環、$V$ を有限 $R$-加群とし、
$J\subset V^{\otimes2}$ とし、$E=T_R(V)/(J)$ を次数一生成の二次 $R$-代数とする。
特に $E_{<0}=0$、$E_0=R$、$E_1=V$ は有限生成である。
$\phi\in E_2$ は中央であると仮定し、$P=E/(\phi)$ と置く。
$P_0,\ldots,P_4$ は有限 projective $R$-加群、$P_5=0$ とする。
さらに、ある $\mathfrak p\in\operatorname{Spec}R$ において
$\phi$ の $E\otimes_RR_{\mathfrak p}$ 上の乗法が単射であると仮定する。
これは剰余体上だけの正則性ではなく、局所環上の正則性の仮定である。

このとき $f\in R\setminus\mathfrak p$ が存在し、単一の開 $D(f)$ 上で
\begin{enumerate}
\item $\phi$ は全次数で左右正則である。
\item 全ての $(E_f)_n$ は有限 projective $R_f$-加群であり、
  \[
  0\longrightarrow(E_f)_{n-2}\xrightarrow{\phi\cdot}(E_f)_n
       \longrightarrow(P_f)_n\longrightarrow0
  \tag{SC-exact}\label{eq:spread-central-regular-sequence}
  \]
  は $R_f$ 上分裂完全である。負次数は零とする。
\item 任意の可換 $R_f$-代数への変更後にも、中央正則性、
  上の全次数の完全列、商 $E_f/(\phi)=P_f$ が保たれる。
\end{enumerate}
さらに $\deg t=2$、$t$ の作用を $\phi$ とすれば、非標準的な次数付き同型
\[
 E_f\simeq\bigoplus_{i=0}^4(P_i)_f\otimes_{R_f}R_f[t](-i)
 \tag{SC-module}\label{eq:spread-central-projective-module}
\]
がある。従って $E_f$ は $R_f[t]$ 上有限 projective でもある。
各 $P_i$ が階数 $1,4,6,4,1$ を持つ場合、$E_n$ の階数は
$1,4,7,8,8,\ldots$ となる。
この補題は $P$ または $E$ の Koszul 性を結論しない。
したがって二次双対 $P^!$ の AS 正則性や全次数平坦性も、この補題だけからは従わない。
\end{lemma}

\begin{proof}
全証明は相対モジュライ論文 \cite[\texttt{noncubic_local_tails}]{QuadricModuli2026} に置く。
\end{proof}

\begin{remark}[中央正則性と Koszul 性の入力]
この補題の有限性は、中央二次元による商が有限個の次数にしか台を持たないことから来る。
したがって $E_n$ ごとに別々の開近傍を選ぶ必要がない。
有限型近傍上で $P_0,\ldots,P_5$ の指定階数と中央性を確立できれば、
局所環で得た $\phi$ の正則性を、全次数について同じ開へ延ばせる。
一方、指定点での $P$ の Koszul 性をその開全体へ延ばすことは別である。
既存の局所 Koszul 変形原理を用いるには、適用先の各局所環の剰余体上の
Koszul 性、またはそれを与える一様な伝播補題がさらに必要である。
第\ref{sec:noncubic-finite-type}節では、全 fiber の有限認識からこの入力を供給し、
補題~\ref{lem:spread-local-central-hull} で同じ根の AS4 被覆まで延長する。
\end{remark}
